\addchap{\lsAcknowledgementTitle}
\markdouble{\lsAcknowledgementTitle}

\noindent Dieses Buch ist eine überarbeitete Fassung meiner Dissertation, die ich im
Dezember 2022 unter dem Titel \textit{Ist} sau cool \textit{cooler als} sehr cool\textit{? Empirische Untersuchungen zur Semantik von deskriptiven und expressiven Ausdrücken am
Beispiel deutscher Intensitätspartikeln} an der Universität des Saarlandes in Saarbrücken eingereicht und im Juli 2023 erfolgreich verteidigt habe. Eine Vielzahl von Menschen hat mir in den vergangenen Jahren mit Rat und Tat zur Seite gestanden und damit wesentlich zum Gelingen meiner Promotion -- und somit auch dieses Buchs -- beigetragen. Allen, die mich auf diesem anstrengenden und mitunter steinigen Weg begleitet und unterstützt haben, gebührt mein aufrichtiger Dank.


Mein großer Dank gilt meinem Betreuer und Doktorvater Prof. Dr. Ingo Reich vom Fachbereich \textit{Germanistik -- Neuere deutsche Sprachwissenschaft} der Universität des Saarlandes für seine beständige Unterstützung und Förderung bei der Entstehung meiner Dissertation. Ihm verdanke ich nicht nur die Möglichkeit zur Promotion am Lehrstuhl, sondern auch mein Interesse an den Themen Expressivität und Intensitätspartikeln. Besonders danke ich ihm für zahlreiche wertvolle Diskussionen über mein Projekt, hilfreiche inhaltliche Anmerkungen und Denkanstöße sowie konstruktive Korrektur- und Optimierungsvorschläge zu früheren Versionen meiner Arbeit.


Ich danke Prof. Dr. Augustin Speyer vom Fachbereich \textit{Germanistik -- Neuere deutsche Sprachwissenschaft} der Universität des Saarlandes, der meine Promotion als Zweitgutachter begleitet hat. Zu großem Dank verpflichtet bin ihm auch für viele inspirierende Anregungen,  nützliche Literaturempfehlungen und bereitwillig ausgeliehene Bücher.


Darüber hinaus bedanke ich mich bei meinen damaligen Kolleg:innen der Abteilung für die fachlichen Ratschläge und Diskussionen, wertvolles Feedback sowie das kritische Prüfen und Korrekturlesen meiner Dissertation. Dank ihnen wird mir meine Promotionszeit immer in guter Erinnerung bleiben. Ein besonderer Dank geht an Robin Lemke und Lisa Schäfer, die ich sowohl in statistischer wie auch methodischer Hinsicht immer wieder aufs Neue um Rat bitten durfte. Ohne ihre unermüdliche Unterstützung und Geduld bei der Auswertung mit R und hilfreichen Anmerkungen zur Statistik wäre diese Arbeit nicht möglich gewesen.


Außerdem danke ich den Produzent:innen der beiden US-amerikanischen Sitcoms \textit{The Big Bang Theory} und \textit{How I Met Your Mother}, durch die ich reichlich fiktive Charaktere und Szenarien für Beispielsätze gewinnen konnte, ohne auf reale Personen zurückgreifen (und diese dadurch schlimmstenfalls unbeabsichtigt diskreditieren) zu müssen.


Ein großes Dankeschön an Richard Page und die Co-Herausgeber:innen Michael Putnam, Laura Catharine Smith, David Natvig und Hanna Fischer für ihr Interesse am Thema sowie für die Aufnahme meiner Arbeit in die Reihe \textit{Open Germanic Linguistics}. Für die sorgfältige Betreuung und wertvolle Unterstützung beim Publikationsprozess möchte ich mich auch bei den Reviewer:innen, den Korrekturleser:innen sowie dem gesamten Team von \textit{Language Science Press} -- insbesondere bei Sebastian Nordhoff -- herzlich bedanken.


Schließlich danke ich meiner Familie und meinen Freund:innen für zahlreiche hilfreiche Tipps zum Durchhalten, aufmunternde Worte und jede weitere Form der Unterstützung, des Verständnisses und des Rückhalts -- sei es in Form von (Nerven‑)Nahrung oder mentaler Ablenkung. Danke, dass ihr all die Jahre an mich geglaubt habt, während aller Höhen und Tiefen für mich da wart und mir stets mit Geduld begegnet seid; selbst dann, wenn ich während intensiver Schreibphasen gereizt, nervlich am Ende oder schlicht wieder einmal nicht verfügbar war. Zu guter Letzt danke ich Lars -- für alles.


